\subsubsection{Primér a štartovací úsek}
Primér je krátky úsek DNA so známou sekvenciou, ktorý sa viaže na komplementárnu časť templátu a určuje miesto, od ktorého sa začína syntéza alebo analýza reťazca.
V biologickom kontexte poskytuje voľnú 3'-OH skupinu pre DNA polymerázu, čím umožňuje predlžovanie nového vlákna.

V tejto práci primér predstavuje zároveň praktický orientačný bod pri spracovaní nanopórových čítaní.
Jeho detekcia umožňuje určiť štartovací úsek reťazca, zjednotiť orientáciu čítaní a odseparovať relevantnú časť sekvencie od technických alebo menej informatívnych oblastí.
Vďaka tomu je možné získať homogennéjšie vstupy pre klastrovanie a následné viacnásobné zarovnanie.

Keďže primér sa nemusí nachádzať presne na začiatku každého čítania, vyhľadáva sa lokálne v celom reťazci.
Tento prístup je robustnejší voči posunu, inzerciám, deléciám a substitúciám, ktoré sú pri nanopórovom sekvenovaní bežné.
